\documentclass{plantillaPPT}

\titulo{13}{Series}

\begin{document}
\primeraDiapo

\begin{frame}{Serie Geométrica}
\framesubtitle{Definición}

\begin{definicion}
    Una serie geométrica es una serie de la forma:
    \[
    \sum_{n=0}^{\infty}ar^n = a+ar+ar^2+ar^3+...
    \]
    Donde \(a\) es el primer término y \(r\) es la razón común.
    \begin{itemize}
        \item Si \(|r|<1\), la serie converge y su valor está dado por \\[1ex]
        \(\displaystyle\sum_{n=0}^{\infty}ar^n = \frac{a}{1-r}\)\\[2ex]
        \item Si \(|r|\geq1\), la serie diverge.
    \end{itemize}
\end{definicion}
    
\end{frame}

\begin{frame}{Serie Geométrica}
\framesubtitle{Sucesión de sumas parciales}

En el caso de una serie geométrica finita de \(n\) términos, la suma de estos viene dada por:
\[S_n=a\cdot\frac{1-r^n}{1-r}\]
    
\end{frame}

\begin{frame}{Modelado de Series}
\framesubtitle{Práctica 13}

1. Considerar un cuadrado de lado 1. Con los puntos medios de sus lados, considerar otro cuadrado y calcular su área. Si el área del cuadrado original es \(a_0\), la del cuadrado interior es \(a_1\), la del cuadrado interior al interior es \(a_2\) y así sucesivamente.\\[3ex]
\((a)\) Calcular las áreas \(a_2,a_3,...,a_n\). \\[2ex]
\((b)\) Expresar la cantidad \(a_0+a_1+a_2+...+a_n\) como una sumatoria.\\[2ex]
\((c)\) ¿La suma anterior es convergente? Justifique su respuesta.\\[2ex]

\end{frame}

\begin{frame}{Valores de Series}
\framesubtitle{Práctica 13}

2. Determinar el valor de la suma:\\[2ex]
\[(a)\;1+\frac{1}{2}+\frac{1}{3}+\frac{1}{4}+\frac{1}{9}+\frac{1}{8}+\frac{1}{27}+\frac{1}{16}+\frac{1}{81}+..._.\]

\[(b)\;\frac{1}{2!}+\frac{2}{3!}+\frac{3}{4!}+\frac{4}{5!}+\frac{5}{6!}+..._.\] \\[2ex]

\textbf{Indicación: }en \((b)\) utilice la definición de serie telescópica.

\end{frame}

\begin{frame}{Valor de Series}
\framesubtitle{Práctica 13}

3. Encontrar la suma de las siguientes series convergentes:
\[(a)\sum_{n=0}^{\infty}\frac{3^{n-1}}{4^{3n+1}} \qquad (b)\sum_{n=5}^{\infty}\frac{3^{n-1}}{4^{3n+1}} \qquad (c)\sum_{n=2}^{\infty}\frac{1}{n^2-1}\]
    
\end{frame}

\begin{frame}{Valor de Series}
\framesubtitle{Práctica 13}
4. Sean \((a_n)_{n\in\mathbb{N}}\) y \((s_n)_{n\in\mathbb{N}}\) sucesiones.\\[1ex]
Si la suma \(a_1+a_2+a_3+...\) es convergente y su sucesión de sumas parciales está dada por \(s_n=\frac{5n^2-3n+2}{n^2-1}\).\\[3ex]

\((a)\) Calcular el valor al que converge la serie \(\sum_{n=1}^{\infty}a_n\).\\[2ex]
\((b)\) Determine el término \(a_n\), considerando la definición de sucesión de sumas parciales.
    
\end{frame}

\end{document}